\begin{helper}
Let \(t\ge2\) be an integer.  There are a natural number \(A=A(t)\) and a
real number \(C=C(t)>0\) such that
\[
  10128t\le A,\qquad A\le C,\qquad 4A\le C,\qquad 2tA\le C\log 2,
\]
and such that every integer \(q\) with \(q\ge C\) satisfies
\[
  q\ge4,\qquad 1\le\log q,\qquad A\le 4q^{t-1},\qquad 1\le Aq^t,
\]
together with the floor-exponent bound
\[
  32tq\log(4q^t)
  <
  \left(\left\lfloor Aq\log q\right\rfloor/(t+1)-1\right),
\]
where the division and subtraction on the right are natural-number truncated
division and subtraction.
\end{helper}
\begin{proof}
Set
\[
  A=
  10128t+100000\,t(t+1)(t+4)+2(t+1)+2
\]
and
\[
  C=\max\left\{4,\ 4A,\ \frac{2tA}{\log2}\right\}.
\]
Since \(\log2>0\), this \(C\) is positive, and the displayed definition gives
\(4\le C\), \(4A\le C\), and \((2tA)/\log2\le C\).  Multiplying the last
inequality by the positive number \(\log2\) gives \(2tA\le C\log2\).  Also
\(A\le4A\le C\), and the first summand in the definition of \(A\) gives
\(10128t\le A\).

Now let \(q\) be an integer with \(q\ge C\).  From \(4\le C\le q\) we get
\(q\ge4\).  Since \(\log3>1\) and the logarithm is increasing on positive
arguments, \(q\ge4\) implies \(1\le\log q\).  Because \(A\le C\le q\) and
\(t\ge2\), we have \(1\le t-1\), so \(q\le q^{t-1}\).  Hence
\[
  A\le q\le q^{t-1}\le4q^{t-1}.
\]
Also \(A\ge1\) and \(q^t\ge1\), so \(1\le Aq^t\).

It remains to prove the floor-exponent estimate.  Put
\[
  B=32tq\log(4q^t).
\]
As \(q\ge4\), the argument \(4q^t\) is at least \(1\), so \(B\ge0\).  We have
\[
  \log(4q^t)=\log4+t\log q.
\]
Moreover \(\log4\le4\le4\log q\), because \(1\le\log q\).  Therefore
\[
  \log(4q^t)\le(t+4)\log q,
\]
and hence
\[
  B\le 32t(t+4)\,q\log q. \tag{1}
\]
Since \(q\log q\ge1\), (1) gives
\[
  (t+1)(B+2)+1
  \le
  \bigl(32t(t+1)(t+4)+2(t+1)+1\bigr)q\log q.
\]
The chosen value of \(A\) is strictly larger than the coefficient
\[
  32t(t+1)(t+4)+2(t+1)+1,
\]
so
\[
  (t+1)(B+2)+1<Aq\log q. \tag{2}
\]

Let \(x=Aq\log q\), \(d=t+1\), and \(n=\lfloor x\rfloor\).  The standard
floor inequality \(x-1<n\), together with (2), gives
\[
  d(B+2)<n.
\]
Dividing by \(d>0\) yields \(B+2<n/d\), where this occurrence of \(n/d\) is
the real quotient.  For natural division, write \(n\div d\) for the quotient.
Then \(n<((n\div d)+1)d\), since the remainder of \(n\) modulo \(d\) is
strictly smaller than \(d\).  Thus the real quotient \(n/d\) is strictly less
than \((n\div d)+1\).  Combining the two inequalities gives
\[
  B+1< n\div d
\]
for the natural quotient.  In particular \(n\div d\ge1\), so natural
subtraction by \(1\) casts to real subtraction, and we conclude
\[
  B< (n\div d)-1.
\]
Substituting \(n=\lfloor Aq\log q\rfloor\) and \(d=t+1\) is exactly the
claimed floor-exponent bound.
\end{proof}
